% ============================================================================
% fpga_secciones.tex  -  Texto y tablas FPGA para el artículo
% Valores: 5 corridas de 60 s en la SP701 (12-nov-2026), potencia Vivado+SAIF.
% Pegar cada bloque en su lugar. Requiere UTF-8 (acentos directos).
% ============================================================================


% ----------------------------------------------------------------------------
% 1) RESUMEN: reemplazar "La caracterización final en FPGA permanece
%    pendiente (TODO)." por:
% ----------------------------------------------------------------------------
En FPGA, un núcleo iterativo RTL validado en hardware procesa 0.047~GB/s a
100~MHz con 1\,751~LUT y 38~mW dinámicos estimados, logrando la menor
energía dinámica neta por byte (0.81~J/GB) a costa de un throughput dos
órdenes de magnitud inferior.

% Sugerencia en la misma frase anterior del resumen: cambiar
%   "la CPU heterogénea alcanza la mayor eficiencia energética del sistema completo"
% por
%   "la CPU heterogénea alcanza la mayor eficiencia energética de sistema completo
%    entre las plataformas con telemetría física"


% ----------------------------------------------------------------------------
% 2) III-A: agregar al final de la sección
% ----------------------------------------------------------------------------
\textit{Adaptación del workload en FPGA:} dado que la BRAM del dispositivo
xc7s100 ($\approx$540~KiB) no puede alojar el archivo de 1~GiB ni el arreglo
de 512~MiB de digests, en FPGA los mensajes se generan en el chip. Cada
mensaje de $C = 64$~B ($B = 2$) se forma con un contador de 64~bits, su
complemento y el digest del mensaje anterior, de modo que cada hash depende
del precedente; el segundo bloque contiene el relleno FIPS~180-4, igual que
en CPU y GPU. Esta modalidad equivale al régimen \textit{kernel-only} de la
GPU, sin transferencia por un bus externo. La correctitud se valida con el
vector NIST ``abc'' ejecutado en silicio tras cada reset y recomputando en el
anfitrión, con \texttt{hashlib.sha256}, la cadena completa de mensajes, cuyo
digest final debe coincidir bit a bit con el de la FPGA.


% ----------------------------------------------------------------------------
% 3) III-B: agregar al final de la sección
% ----------------------------------------------------------------------------
En FPGA, el diseño opera a 100~MHz, obtenidos dividiendo por dos el reloj
diferencial de 200~MHz de la placa. Un núcleo \textit{Virtual I/O} (VIO)
controla el inicio y la parada del banco de prueba y expone el número de
mensajes completados, los ciclos activos y el último digest. Un script Tcl en
el \textit{Hardware Manager} de Vivado aplica el mismo protocolo que en CPU y
GPU (3~s de calentamiento y 5~corridas, aquí de 60~s), cronometradas con el
reloj del anfitrión. El throughput (Ec.~3) se calcula con los ciclos contados
en hardware, y el reloj del anfitrión se usa para verificar la frecuencia.
La potencia se estima con \texttt{report\_power} de Vivado~2026.1 sobre el
diseño post-ruteo sin núcleos de depuración, con actividad de conmutación
obtenida de una simulación funcional post-implementación de 1~ms en régimen
continuo (SAIF; 58\,\% de las nets anotadas y el resto por propagación
probabilística; confianza \textit{High}). Como aproximación a la Ec.~4, la
energía neta usa la potencia dinámica del bloque evaluado y la energía bruta
suma la potencia estática del dispositivo (0.089~W).


% ----------------------------------------------------------------------------
% 4) III-E: agregar como limitación (9) al final de la lista
% ----------------------------------------------------------------------------
y (9) en FPGA la potencia es una estimación del modelo de Vivado con
actividad simulada y no una medición física como RAPL o NVML, y abarca solo
el chip (NVML incluye además la memoria y la placa de la GPU); asimismo, el
workload se genera en el chip y no incluye transferencias desde el
anfitrión, por lo que la FPGA se compara con el régimen \textit{kernel-only}
de la GPU y no con las mediciones E2E.


% ----------------------------------------------------------------------------
% 5) III-F. Diseño de Microarquitectura RTL  (texto de la sección vacía)
% ----------------------------------------------------------------------------
El núcleo (\texttt{sha256\_core}) implementa una arquitectura iterativa de
una ronda por ciclo (\textit{loop rolling} [6]). Una máquina de estados de
tres fases carga el bloque de 512~bits y las variables $a$--$h$ (1~ciclo),
ejecuta las 64~rondas (64~ciclos) y suma el estado intermedio
$H^{(i)} = H^{(i-1)} + \{a,\dots,h\}$ (1~ciclo), totalizando 66~ciclos por
bloque. Las constantes $K_t$ se implementan como una ROM en LUT, y el
encadenamiento de $H^{(i)}$ permite procesar mensajes de $B \geq 2$ bloques.
Un banco de prueba en el chip (\texttt{sha256\_bench}) alimenta el núcleo con
mensajes de $C = 64$~B y añade 2~ciclos de intercambio (\textit{handshake})
por bloque, por lo que la latencia medida es de 68~ciclos por bloque y 136
por mensaje.


% ----------------------------------------------------------------------------
% 6) III-G. Optimizaciones de Ruta de Datos y Ruta de Control
% ----------------------------------------------------------------------------
La expansión de mensajes usa una ventana deslizante de 16~palabras
($W[0..15]$, 512~FF) en lugar de almacenar las 64~palabras $W_t$ (2\,048~FF),
técnica estándar [6] análoga a la ventana \texttt{w[16]} del kernel GPU~v2.
En cada ronda, la nueva palabra
$W_{t+16} = \sigma_1(W_{t+14}) + W_{t+9} + \sigma_0(W_{t+1}) + W_t$ se
calcula en paralelo y fuera del camino crítico del compresor, que queda
determinado por la suma $T_1 + T_2$. La síntesis mapea los sumadores sobre
cadenas de acarreo CARRY4 y no utiliza bloques DSP ni BRAM; además, Vivado
recodificó la máquina de estados en \textit{one-hot} y replicó sus registros
por \textit{fanout}, lo que añade 20~FF a los 1\,034~FF del RTL. No se
aplicaron técnicas de reducción del camino crítico, como el precómputo de
$h + K_t + W_t$ [8], el reordenamiento espacial o el
\textit{quasi-pipelining} [6]; su incorporación queda como trabajo futuro.


% ----------------------------------------------------------------------------
% 7) TABLA IV: reemplazar la fila "FPGA (Spartan-7) TODO" por:
% ----------------------------------------------------------------------------
FPGA Spartan-7 (1 núcleo)$^{b}$ & 0.0471 $\pm$ 0.0000 & 0.7353 $\pm$ 0.0000
  & 0.81$^{c}$ & --- & 0.013$\times^{b}$ \\

% Notas a agregar bajo la Tabla IV:
% $^{b}$ Régimen en chip, comparable con GPU~v2 kernel y no con E2E.
%   Throughput calculado con los ciclos contados en hardware a 100~MHz
%   (variación entre corridas $<0.01$\,\%).
% $^{c}$ Potencia dinámica estimada (Vivado, SAIF) del banco completo
%   (0.038~W); núcleo solo: 0.72~J/GB. Incluyendo la potencia estática del
%   dispositivo (0.089~W): 2.70~J/GB. Sin anfitrión activo durante la medición.


% ----------------------------------------------------------------------------
% 8) TABLA V: reemplazar la tabla completa
% ----------------------------------------------------------------------------
\begin{table}[!t]
\caption{Resultados del acelerador RTL SHA-256 en FPGA (AMD Spartan-7)}
\label{tab:fpga}
\centering
\footnotesize
\begin{tabular}{@{}ll@{}}
\toprule
Parámetro & Valor \\
\midrule
Dispositivo / placa              & xc7s100fgga676-2 / SP701 \\
HDL / arquitectura               & Verilog RTL / 1 ronda por ciclo \\
Frecuencia de operación          & 100~MHz \\
$f_{\text{máx}}$ estimada (STA)$^{a}$ & $\geq$117.2~MHz ($\geq$121.3 sin VIO) \\
Ciclos / bloque                  & 66 (núcleo) / 68 (banco, medido) \\
Ciclos / mensaje de 64~B         & 136 (medido) \\
Mbps por bloque                  & 775.76 (núcleo) / 752.94 (banco) \\
Throughput útil ($C = 64$~B)     & 0.0471~GB/s (0.7353~Mhash/s) \\
LUT                              & 1\,751 / 64\,000 (2.74\,\%) \\
FF$^{b}$                         & 1\,054 / 128\,000 (0.82\,\%) \\
BRAM / DSP                       & 0 / 0 \\
Eficiencia de área               & 0.443~Mbps/LUT \\
Potencia dinámica (SAIF)$^{c}$   & 0.034 (núcleo) / 0.038~W (banco) \\
Potencia estática (chip)         & 0.089~W \\
Energía neta                     & 0.72 (núcleo) / 0.81~J/GB (banco) \\
Energía bruta                    & 2.61 (núcleo) / 2.70~J/GB (banco) \\
Speedup vs.\ CPU SHA-NI 4T       & 0.013$\times$ \\
\bottomrule
\multicolumn{2}{@{}p{0.95\columnwidth}@{}}{\scriptsize
$^{a}$ Holgura positiva con restricción de 10~ns (WNS = 1.468~ns con VIO,
1.759~ns sin VIO); cota inferior, sin barrido de restricción.
$^{b}$ 1\,034~FF del RTL más 20 añadidos por síntesis (FSM \textit{one-hot}
replicada). $^{c}$ Estimación post-ruteo con actividad de simulación
post-implementación (confianza \textit{High}). Recursos del núcleo sin la
infraestructura de depuración (VIO y \textit{dbg\_hub}).}
\end{tabular}
\end{table}

% Si la plantilla no carga booktabs, cambiar \toprule, \midrule y \bottomrule
% por \hline.


% ----------------------------------------------------------------------------
% 9) IV-D. Resultados de FPGA: reemplazar el texto actual (incluido el
%    párrafo "La implementación FPGA se enfoca en maximizar...")
% ----------------------------------------------------------------------------
La Tabla~\ref{tab:fpga} resume los resultados post-implementación y en
hardware del acelerador sobre la placa SP701. En 5~corridas de 60~s a
100~MHz, el banco completó $44.14 \pm 0.005$~millones de mensajes por
corrida ($\approx$2.8~GB útiles) con exactamente 136~ciclos por mensaje en
todas ellas. El digest final de cada corrida coincidió bit a bit con la
cadena recomputada con \texttt{hashlib.sha256} (231.8~millones de mensajes
encadenados en total), y el autotest NIST pasó en silicio. Esto equivale a
0.7353~Mhash/s y 0.0471~GB/s útiles (752.94~Mbps por bloque). La frecuencia
estimada con el reloj del anfitrión fue $99.976 \pm 0.004$~MHz; la diferencia
de 0.02\,\% respecto del valor nominal corresponde a la latencia JTAG de los
comandos de inicio y parada (11--17~ms por corrida). El diseño cumple las
restricciones de temporización con una holgura de 1.468~ns sobre 10~ns
(1.759~ns sin núcleos de depuración), lo que indica una $f_{\text{máx}}$ de
al menos 117~MHz.

Excluyendo la infraestructura de depuración (VIO y \textit{dbg\_hub}:
1\,659~LUT y 2\,936~FF), el núcleo ocupa 1\,751~LUT (2.74\,\%) y 1\,054~FF
(0.82\,\%), sin DSP ni BRAM, con una eficiencia de área de 0.443~Mbps/LUT.
Con actividad de simulación (SAIF), la potencia dinámica estimada es de
34~mW para el núcleo y 38~mW para el banco completo, más del doble de la
estimación sin simulación (16 y 22~mW), lo que indica que el modo
\textit{vectorless} subestima la conmutación de SHA-256. Esto resulta en
0.72~J/GB netos para el núcleo y 0.81~J/GB para el banco, por debajo de la
GPU~v2 kernel-only (1.62~J/GB) y de la CPU con 12~CPUs (P+E) (2.53~J/GB).
Al sumar la potencia estática del dispositivo (0.089~W), la energía bruta
del banco es de 2.70~J/GB: inferior a la de la CPU ($\approx$3.41~J/GB con
su potencia de reposo) pero superior a la de la GPU~v2 kernel-only
($\approx$1.90~J/GB), porque con un solo núcleo la potencia estática
representa el 70\,\% del total del chip. Con 0.127~W en el chip, la FPGA
opera dos órdenes de magnitud por debajo de la CPU (20~W) y la GPU, a costa
de un throughput 124~veces menor que la mejor configuración de CPU.


% ----------------------------------------------------------------------------
% 10) IV-C: en la última frase del punto "Eficiencia energética de
%     dispositivo vs. sistema completo", cambiar
%     "consolidando a la CPU heterogénea como la arquitectura más eficiente"
%     por
%     "consolidando a la CPU heterogénea como la arquitectura más eficiente
%      entre CPU y GPU"
% ----------------------------------------------------------------------------


% ----------------------------------------------------------------------------
% 11) V. CONCLUSIONES: reemplazar desde "Como trabajo futuro, la
%     caracterización final del acelerador RTL..." hasta el final por:
% ----------------------------------------------------------------------------
La caracterización en FPGA completa la comparativa de tres vías: un núcleo
iterativo RTL en Spartan-7, validado en hardware sobre 231.8~millones de
mensajes encadenados, procesa 0.047~GB/s a 100~MHz con 1\,751~LUT y
0.127~W en el chip. Alcanza la menor energía dinámica neta estimada
(0.81~J/GB), pero su throughput es dos órdenes de magnitud inferior al de
CPU y GPU, y su energía bruta (2.70~J/GB) está dominada por la potencia
estática del dispositivo. Como trabajo futuro se plantea replicar el núcleo
y acortar su camino crítico para amortizar la potencia estática, medir la
potencia directamente en la placa y evaluar la FPGA de extremo a extremo con
una interfaz de alta velocidad.
